\documentclass[11pt,a4paper]{article}
\usepackage[margin=2cm]{geometry}
\usepackage[T1]{fontenc}
\usepackage[utf8]{inputenc}
\usepackage{microtype}
\usepackage{setspace}
\usepackage{parskip}
\usepackage{titlesec}
\usepackage{fancyhdr}
\usepackage{booktabs}
\usepackage{tabularx}
\usepackage{xcolor}
\usepackage{listings}
\usepackage{amsmath}
\usepackage{hyperref}
\usepackage{enumitem}
\usepackage{tcolorbox}
\usepackage{multicol}

\definecolor{amber}{RGB}{245,158,11}
\definecolor{zinc}{RGB}{63,63,70}

\onehalfspacing
\setlength{\parindent}{0pt}
\setlength{\parskip}{5pt}
\titleformat{\section}{\normalsize\bfseries}{{\color{amber}\thesection.}}{0.4em}{}[\vspace{-2pt}\rule{\linewidth}{0.3pt}]
\titleformat{\subsection}{\small\bfseries}{\thesubsection}{0.4em}{}

\pagestyle{fancy}
\fancyhf{}
\fancyhead[L]{\small\textcolor{zinc}{Smart Alloy Selector}}
\fancyhead[R]{\small\textcolor{zinc}{MET-QUEST'26}}
\fancyfoot[C]{\small\thepage}
\renewcommand{\headrulewidth}{0.3pt}

\tcbuselibrary{skins}
\newtcolorbox{highlight}{
  colback=amber!8,colframe=amber,boxrule=0.5pt,arc=3pt,
  left=5pt,right=5pt,top=3pt,bottom=3pt}

\hypersetup{colorlinks,linkcolor=amber!80!black,urlcolor=blue}

\begin{document}

% -- TITLE --------------------------------------------------
\begin{center}
{\Large\bfseries Smart Alloy Selector}\\[4pt]
{\color{amber}\large Needle in the Data-Stack --- MET-QUEST'26}\\[4pt]
{\small\textbf{Live:} \url{https://beautiful-tiramisu-211a94.netlify.app}
\quad \textbf{Code:} \url{https://github.com/akshatm24/mse_hackathon}}\\[4pt]
\begin{tabular}{ll|ll}
  \textbf{LLM} & Gemini Flash 2.0 &
  \textbf{Materials} & 150+ (MP API + curated) \\
  \textbf{Stack} & Next.js 14 + TypeScript &
  \textbf{Deploy} & Netlify serverless \\
\end{tabular}
\end{center}

\vspace{-4pt}\rule{\linewidth}{0.5pt}\vspace{4pt}

% -- ABSTRACT -----------------------------------------------
\begin{highlight}
\textbf{Abstract.}
Smart Alloy Selector is an AI-powered web tool that replaces
manual database lookup in material selection. A user describes
an engineering problem in plain English; the system uses Gemini
Flash 2.0 to extract quantitative constraints, runs a validated
weighted scoring engine against 150+ materials sourced from
The Materials Project API, ASM Handbook, and manufacturer
datasheets, and returns ranked candidates with a natural-language
explanation. The tool degrades gracefully when the LLM is
unavailable using a keyword-based heuristic extractor.
\end{highlight}

% -- SECTION 1 ----------------------------------------------
\section{Introduction and Problem}

Material selection is a critical engineering bottleneck.
While databases like MatWeb and ASM Handbook contain extensive
data, they are static, require expert interpretation, and give
no ranking or trade-off analysis for a given design requirement.
An engineer building a 4-point probe station for sintered
copper-cobalt pellets must simultaneously evaluate electrical
conductivity, hardness, thermal stability to 200$^\circ$C, and chemical
inertness --- a search across multiple disconnected tables.

MET-QUEST'26 requires an AI tool acting as a virtual materials
scientist: given user-defined constraints, query a property
database and recommend the optimal material with a clear
explanation. The Brownie Point requirement adds a natural-language
interface and property prediction for novel compositions.

\textbf{Our approach:} natural language in $\rightarrow$ AI extracts
constraints $\rightarrow$ deterministic scoring $\rightarrow$ AI explains results.
The two LLM calls (extraction + explanation) are isolated
from the scoring engine, so the tool is fully usable even
without a working API key.

% -- SECTION 2 ----------------------------------------------
\section{System Architecture}

\subsection{Pipeline}

The system operates as three sequential stages per request:

\textbf{Stage 1 --- Constraint Extraction.}
The query is sent to Gemini Flash 2.0 with a system prompt
requiring a JSON object of type \texttt{UserConstraints}.
Critical rule enforced in the prompt: if the user says
``cheap'', the cost weight must be $\geq 0.60$; if they say
``lightweight'', the density weight must be $\geq 0.60$.
On quota error (HTTP 429), a local heuristic extractor runs
instead, counting keyword hits per property axis and setting
weights proportionally.

\textbf{Stage 2 --- Deterministic Scoring.}
Category intent is inferred first (e.g. FDM query $\rightarrow$ exclude
ceramics). Hard filters eliminate candidates violating strict
constraints. Surviving candidates are scored with:
\begin{equation}
s_i = 100 \cdot \!\sum_{k} w_k \hat{p}_{ik}, \quad
\hat{p}_{ik} \in [0,1], \quad \textstyle\sum_k w_k = 1
\end{equation}
where $\hat{p}_{ik}$ is the min-max normalised property $k$
of material $i$, computed \emph{within the filtered set only}
(not the full database), and cost uses log-scale normalisation
to prevent outliers like AuSn20 (\$40{,}000/kg) from collapsing
the cost axis.

\textbf{Stage 3 --- Explanation.}
Top-5 materials + original query are sent to Gemini for a
3-paragraph prose explanation. A template-based fallback
activates on API failure.

\subsection{Technology Stack}

\begin{tabular}{lll}
\toprule
\textbf{Layer} & \textbf{Tech} & \textbf{Why} \\
\midrule
Frontend & Next.js 14, App Router &
  Server-side routes hide the API key \\
LLM & Gemini Flash 2.0 &
  Free tier, $<$1s, reliable JSON mode \\
Charts & Recharts &
  Built-in RadarChart, Next.js SSR compatible \\
Deploy & Netlify &
  Auto-deploy from GitHub, serverless functions \\
\bottomrule
\end{tabular}

% -- SECTION 3 ----------------------------------------------
\section{Materials Database (150+ entries)}

\subsection{Data Sources and Access Procedure}

\begin{description}[leftmargin=0pt,itemsep=2pt]
  \item[\textbf{The Materials Project API} (automated, primary)]
  \url{materialsproject.org} provides DFT-computed properties
  for 154{,}000+ materials via a free REST API.
  Access: register at \url{next.materialsproject.org/api},
  copy the API key to \texttt{.env.local} as
  \texttt{MATERIALS\_PROJECT\_API\_KEY}, then run
  \texttt{npm run update-db}. The script fetches density,
  elastic moduli (bulk + shear $\rightarrow$ Young's via VRH average),
  band gap, and is\_metal flag, then derives corrosion
  resistance and service temperature from lookup tables.

  \item[\textbf{ASM Handbook} (hand-curated, metals)]
  Volumes 1 and 2 provide peer-reviewed room-temperature
  properties for all steels, aluminium alloys, and
  superalloys in the curated set.
  Access: \url{asminternational.org} (institutional
  subscription or purchase). Properties extracted: tensile
  strength, yield strength, hardness, thermal conductivity,
  thermal expansion, melting point.

  \item[\textbf{MatWeb} (hand-curated, polymers + ceramics)]
  \url{matweb.com} free tier provides property tables for
  most thermoplastics and engineering ceramics.
  Access: search by material name $\rightarrow$ copy property table.
  Premium tier offers API access (REST, JSON) with a
  subscription key.

  \item[\textbf{Manufacturer Datasheets} (hand-curated)]
  Used for specialty materials: Victrex (PEEK), Stratasys
  (Ultem), Special Metals (Inconel, Waspaloy, Monel),
  Solvay (PVDF, PPS), Arkema (PEKK), Indium Corporation
  (solder alloys). All PDFs freely available on manufacturer
  websites.

  \item[\textbf{NASA TPSX} (hand-curated, ceramics)]
  \url{tpsx.arc.nasa.gov} (free, requires NASA account).
  Used for aerospace ceramics and carbon-carbon composite.
  Properties: high-temperature thermal conductivity and
  specific heat as CSV curves.
\end{description}

\subsection{Database Summary}

\begin{tabular}{lrp{7cm}}
\toprule
\textbf{Category} & \textbf{Count} & \textbf{Examples} \\
\midrule
Metals \& Alloys & 60+ &
  316L SS, Ti-6Al-4V, Inconel 718, Al 6061, Cu C110,
  pure metals (Fe, Ni, Co, Pt, Au, Ag), Invar, Kovar \\
Polymers & 25+ &
  PEEK, PETG, Nylon PA12, PEKK, Ultem, PPS, TPU,
  Polyimide, Silicone, PVC, PS, PP \\
Ceramics & 20+ &
  Al$_2$O$_3$, ZrO$_2$, SiC, Si$_3$N$_4$, AlN,
  TiB$_2$, ZrB$_2$, WC-Co, TiO$_2$ \\
Composites & 8+ &
  CFRP, GFRP, Kevlar/Epoxy, C-C, SiC/SiC CMC \\
Solders & 8+ &
  SAC305, Sn63Pb37, AuSn20, SnBi42, InBi57, BAg-7 \\
\midrule
\textbf{Total} & \textbf{150+} & \\
\bottomrule
\end{tabular}

\subsection{Properties Stored (18 fields per material)}

Density, tensile strength, yield strength, elastic modulus,
Vickers hardness, thermal conductivity, specific heat,
melting point, glass transition temperature,
max service temperature, thermal expansion (CTE),
electrical resistivity, corrosion resistance (ordinal),
machinability (ordinal), FDM printability (ordinal),
cost USD/kg (2024 approximate), application tags, data source.

% -- SECTION 4 ----------------------------------------------
\section{Scoring Methodology}

\subsection{Stage 1: Category Intent Filter}
Before numeric scoring, the raw query is scanned for
category-level signals. This prevents ceramics from appearing
in FDM plastic queries regardless of their extreme property
values.

\begin{tabular}{ll}
\toprule
\textbf{Query signal} & \textbf{Action} \\
\midrule
3d print / fdm / filament & Exclude Ceramic, Solder \\
solder / bga / reflow     & Include only Solder, Metal \\
polymer / plastic         & Exclude Ceramic, Solder \\
No signal                 & No category exclusion \\
\bottomrule
\end{tabular}

\subsection{Stage 2: Hard Filtering}
Eliminate if: $T_{max}^{(i)} < c_T$ \ OR\
$\sigma_t^{(i)} < c_\sigma$ \ OR\
$\rho^{(i)} > c_\rho$ \ OR\ $C^{(i)} > c_C$ \ OR\
corrosion rank $< $ required rank, OR FDM/conductivity
binary constraints violated. Four progressive fallback passes
relax numeric and category constraints until $\geq 3$
candidates survive.

\subsection{Stage 3: Weighted Scoring}
\begin{equation}
s_i = 100 \cdot \bigl(
  w_T\hat{T}_i + w_S\hat{\sigma}_i +
  w_W(1-\hat{\rho}_i) + w_C(1-\hat{C}_i^{\log}) +
  w_R\hat{R}_i \bigr), \quad s_i \in [0,100]
\end{equation}
Cost uses $\hat{C}_i^{\log} = \text{norm}(\ln(1+C_i))$
to handle the \$1--\$62{,}000/kg range. All normalisations
use min-max within the filtered set only. Weights are
re-normalised to sum~=~1.0 before scoring.

\subsection{Weight Inference}
\begin{tabular}{lrrrrr}
\toprule
\textbf{Dominant signal} & $w_T$ & $w_S$ & $w_W$ & $w_C$ & $w_R$ \\
\midrule
``cheap'' / ``budget''      & .07 & .13 & .07 & \textbf{.65} & .08 \\
``lightweight'' / ``drone'' & .13 & .13 & \textbf{.65} & .07 & .02 \\
``strong'' / ``structural'' & .13 & \textbf{.65} & .10 & .07 & .05 \\
``heat'' / ``motor''        & \textbf{.65} & .13 & .10 & .07 & .05 \\
``marine'' / ``acid''       & .10 & .15 & .05 & .05 & \textbf{.65} \\
No signal (default)         & .18 & .28 & .18 & .28 & .08 \\
\bottomrule
\end{tabular}

% -- SECTION 5 ----------------------------------------------
\section{Worked Examples}

\textbf{Example 1 --- 3D-printed motor bracket (problem statement)}\\
Query: \textit{``lightweight bracket for DC motor, 3D printed,
cannot warp at 85$^\circ$C''}\\
Inferred: $c_T$=85$^\circ$C, FDM=true, $w_T$=0.40, $w_W$=0.35\\
Filtered set: 9 FDM polymers survive. Winner: \textbf{PETG}
(score 88/100) --- exactly meets thermal limit, lowest density
among survivors, excellent desktop printability.

\textbf{Example 2 --- 4-point probe for Cu-Co pellets}\\
Query: \textit{``probe tips: electrically conductive, hard,
stable to 200$^\circ$C''}\\
Inferred: conductive=true, $c_T$=200$^\circ$C, $w_S$=0.40, $w_T$=0.30\\
Filtered: 14 metals (all insulators eliminated).
Winner: \textbf{Beryllium Copper C17200}
($\sigma_t$=1380 MPa, $\rho_e$=7$\times10^{-8}$ $\Omega$·m).

\textbf{Example 3 --- Cheapest structural material}\\
Query: \textit{``cheapest material for a bracket, no special needs''}\\
Inferred: $w_C$=0.65 (cost dominant), no hard constraints\\
Winner: \textbf{Grey Cast Iron} or \textbf{Carbon Steel 1020}
(\$0.5--\$0.8/kg) --- cost efficiency score dominates.

% -- SECTION 6 ----------------------------------------------
\section{Security, Deployment, and Limitations}

\textbf{Security.}
The Gemini API key is stored only in \texttt{.env.local}
(listed in \texttt{.gitignore}) and in Netlify's encrypted
environment variable store. It is read only inside the
serverless function via \texttt{process.env.GEMINI\_API\_KEY}
and never exposed to the browser bundle. The Materials Project
key is handled identically. The \texttt{.env.example} committed
to the public repo contains placeholder text only.

\textbf{Deployment.}
Auto-deploys on every \texttt{git push} to \texttt{main}.
Build time: $\approx$90~seconds. The \texttt{GEMINI\_API\_KEY}
and \texttt{MATERIALS\_PROJECT\_API\_KEY} variables are set in
Netlify Dashboard $\rightarrow$ Site Settings $\rightarrow$ Environment Variables.

\textbf{Limitations.}
All properties are at room temperature. No fatigue life,
fracture toughness, or property-temperature curves. Cost
values are 2024 market approximations. LLM may misinterpret
highly ambiguous queries; inferred weights are displayed
to users for transparency.

\textbf{Future work.}
Property-temperature curves from ASM data; ML composition
interpolation using AFLOW/OQMD (3.5M+ DFT entries);
fatigue data from MMPDS; live supplier pricing API.

% -- REFERENCES ---------------------------------------------
\begin{thebibliography}{9}
\bibitem{mp} A. Jain et al., ``Commentary: The Materials Project,''
  \textit{APL Materials}, 1(1):011002, 2013.
\bibitem{asm} ASM International,
  \textit{ASM Handbook Vol 1 \& 2}, 10th ed., 2023.
\bibitem{matweb} MatWeb LLC,
  \textit{Material Property Data}, \url{matweb.com}, 2024.
\bibitem{ashby} M. F. Ashby,
  \textit{Materials Selection in Mechanical Design},
  4th ed. Butterworth-Heinemann, 2011.
\bibitem{gemini} Google DeepMind,
  \textit{Gemini Technical Report}, 2024.
\bibitem{nasa} NASA Ames,
  \textit{TPSX Database}, \url{tpsx.arc.nasa.gov}, 2023.
\end{thebibliography}

\end{document}

% -- HOW TO COMPILE -----------------------------------------
% Option 1 (easiest): Upload to overleaf.com -> Compile -> PDF
% Option 2 (local):   pdflatex smart_alloy_report.tex
